% Condensed version of muon/notes_spectral_theory.tex: only the main findings, in the format of a short paper.
% Build (from muon/):  pdflatex paper_spectral_muon && bibtex paper_spectral_muon && pdflatex paper_spectral_muon && pdflatex paper_spectral_muon
\documentclass[11pt]{article}
\usepackage[utf8]{inputenc}
\usepackage[T1]{fontenc}
\usepackage{amsmath,amssymb,amsthm,mathtools}
\usepackage{graphicx}
\usepackage{booktabs}
\usepackage{enumitem}
\usepackage[margin=1in]{geometry}
\usepackage[colorlinks=true,linkcolor=blue!55!black,citecolor=blue!55!black,urlcolor=blue!55!black]{hyperref}
\usepackage{xcolor}
\graphicspath{{figures/}}
\setlength{\emergencystretch}{2em}

\newtheorem{theorem}{Theorem}
\newtheorem{proposition}[theorem]{Proposition}
\newtheorem{lemma}[theorem]{Lemma}
\theoremstyle{remark}
\newtheorem{remark}[theorem]{Remark}

\newcommand{\R}{\mathbb{R}}
\newcommand{\E}{\mathbb{E}}
\newcommand{\Prob}{\mathbb{P}}
\newcommand{\norm}[1]{\left\|#1\right\|}
\newcommand{\nop}[1]{\left\|#1\right\|_{\mathrm{op}}}
\newcommand{\nnuc}[1]{\left\|#1\right\|_{*}}
\newcommand{\nF}[1]{\left\|#1\right\|_{F}}
\newcommand{\inner}[2]{\left\langle #1, #2 \right\rangle}
\newcommand{\abs}[1]{\left|#1\right|}
\newcommand{\tr}{\operatorname{tr}}
\newcommand{\diag}{\operatorname{diag}}
\newcommand{\polar}{\operatorname{polar}}
\newcommand{\srank}{\operatorname{srank}}
\newcommand{\nr}{\operatorname{nr}}
\newcommand{\Var}{\operatorname{Var}}
\newcommand{\MAD}{\operatorname{MAD}}
\newcommand{\Med}{\operatorname{Med}_\ast}
\newcommand{\Lop}{L_{\mathrm{op}}}
\newcommand{\muop}{\mu_{\mathrm{op}}}
\newcommand{\kop}{\kappa_{\mathrm{op}}}
\newcommand{\kF}{\kappa_{F}}
\newcommand{\cO}{\mathcal{O}}
\newcommand{\lmin}{\lambda_{\min}}
\newcommand{\lmax}{\lambda_{\max}}
\newcommand{\ones}{\mathbf{1}}

\title{What limits Muon on strongly convex problems?\\ Conditioning barriers and convergence to a nuclear median}
\author{Side study of the project ``Accelerating Variance Reduction''\\ Statistical Learning Theory, UC Berkeley}
\date{16 November 2026}

\begin{document}
\maketitle

\begin{abstract}
Muon updates a matrix parameter along the polar factor of its (averaged) gradient: steepest descent for the spectral
norm. We ask two questions inherited from a study of variance-reduced methods: can this geometry improve worst-case
constants, and what does gradient noise do to the limit of the iteration? On strongly convex problems we obtain three
groups of results. (i) \emph{Conditioning.} For ridge-regularized generalized linear models whose loss curvature
vanishes along a sequence of parameters (logistic, softmax), no norm, spectral or other, has a better worst-case
condition number than the Euclidean one. In the spectral geometry the condition number is at least $\min(K,d)$, a
folklore barrier; we add that equality holds only for isotropic quadratics, a frame bound and a sharp comparison with
the Euclidean condition number. (ii) \emph{Noise.} With additive gradient noise, stochastic polar steps are stochastic
mirror descent, in the space of gradients, on a noise-smoothed nuclear norm. For step lengths with
$\sum_t\eta_t=\infty$ and $\sum_t\eta_t^2<\infty$ they converge almost surely to a point where zero is a \emph{nuclear
median}, rather than the mean, of the stochastic gradient, and the set of possible limits does not depend on the
schedule. For batch size one on least squares the limit is a weighted least-distances estimator; for skewed gradient
noise with independent entries a formal expansion gives an offset of order $1/b$, whose constant we compute numerically
and check by simulation: per entry it is smaller than for sign descent by a factor of order $1/d$; a finite-sum
variance-reduced estimator has a fixed point without offset. (iii) \emph{Dynamics.} On multi-output least squares,
whose polar dynamics are largely known, we add a conservation law in dimension two, a counterexample showing that the
spectral norm of the error can increase along the polar flow in dimension three, a sharp contraction bound and bounds
on extinction times. Apart from the published facts recalled as such, these statements were not found in the
literature checks we ran; a few refinements obtained during the review were not searched separately and are listed as
such.
\end{abstract}

\section{Introduction}

Muon \cite{jordan2024muon} replaces the gradient $G$ of each weight matrix by an approximation of its polar factor
$\polar(G)=UV^\top$ ($G=U\diag(s)V^\top$): it is steepest descent for the spectral norm
\cite{carlson2015stochastic,bernstein2024old}, a linear minimization oracle over a spectral ball
\cite{pethick2025training}. A rapidly growing literature analyses this step on models where everything can be computed:
matrix quadratics \cite{shulgin2026quadratic,jawanpuria2026constant,gonon2026insights,davis2025spectral,wang2026why},
linear models and factorizations \cite{vasudeva2025muon,ma2026preconditioning}, and noisy gradients
\cite{xie2026muon,chen2026when,paquette2026phases,do2026muon}.

This note comes from a different direction. Our main project studies variance-reduced stochastic methods for
finite sums (SAGA \cite{defazio2014saga}, SVRG \cite{johnson2013accelerating}, the analysis of \cite{bach2024learning}),
and reached two conclusions: worst-case constants hide what happens along trajectories, and no fixed quadratic metric
improves the worst-case condition number of regularized logistic regression. We asked what these conclusions become for
Muon, and what a variance-reduction viewpoint says about its \emph{speed} and its \emph{precision}. A first set of
answers was submitted to an adversarial review (three passes on the mathematics, three on the literature); a large part
turned out to be published in 2025--2026, and is cited as such below. What remains is the following.

\paragraph{Contributions.}
\begin{itemize}[leftmargin=1.3em,itemsep=1pt,topsep=2pt]
\item \textbf{No free lunch for every norm} (Theorem~\ref{thm:nfl}). For ridge-regularized generalized linear models
whose loss curvature vanishes along a sequence of parameters (logistic, softmax), every norm $N$ satisfies
$\kappa_N\ge\kF$, for the full objective and for the per-sample constants that govern variance-reduced methods. The
spectral condition number obeys $\kop\ge\min(K,d)$ (folklore), which we refine by a frame bound, a rigidity statement
and a sharp two-sided comparison with $\kF$ (Proposition~\ref{prop:frame}).
\item \textbf{Convergence to a nuclear median} (Theorem~\ref{thm:smd}). With additive noise $\xi$, the iteration
$W_{t+1}=W_t-\eta_t\polar(\nabla F(W_t)+\xi_t)$ is stochastic mirror descent on $\psi_\xi(G)=\E[\nnuc{G+\xi}-\nnuc\xi]$ with
mirror map $F^\ast$. Hence, for lengths with $\sum_t\eta_t=\infty$ and $\sum_t\eta_t^2<\infty$, $\nabla F(W_t)$
converges almost surely to a minimizer of $\psi_\xi$, the negative of a nuclear median of $\xi$: if $0$ is not such a
median, $F(W_t)-F_\ast$ has a positive limit, with the same lower bound for every schedule of this class. Companion
statements: a least-distances limit for batch one on least squares (Proposition~\ref{prop:lad}), a formal first-order
offset of order $1/b$ for skewed noise with independent entries, with a dimension-dependent constant
(Proposition~\ref{prop:offset}), and a fixed point without offset for finite-sum variance-reduced estimators
(Proposition~\ref{prop:median}(d)).
\item \textbf{Complements on the least-squares dynamics} (Propositions~\ref{prop:Q}--\ref{prop:ext}): an exact
conservation law for $2\times2$ problems, an integer counterexample to the monotonicity of $\nop{W_t-W_\ast}$ along the
polar flow, the bound $1-2/\kop$ on the contraction of the nuclear-scaled step after its first iteration in the
commuting case, and bounds on the first rank-drop time and, for $d=2$, on the extinction time of the flow.
\end{itemize}

\paragraph{Related work.} \emph{Geometry and rates.} Smoothness constants for general norms and the rate $1-\mu/L$
of steepest descent under a dual-norm Polyak--\L{}ojasiewicz inequality are in \cite{balles2020geometry}; spectral-norm
versions for Muon in \cite{riabinin2025gluon,shen2025convergence,kovalev2025noneuclidean} (in
\cite{lau2025polargrad} with Frobenius constants, hence an extra rank factor); the constant
$\tr\Sigma$ and the stable-rank criterion for least-squares layers in \cite{davis2025spectral}. That no function is
well conditioned for $\ell_\infty$ in high dimension is folklore \cite{sidford2018coordinate}. \emph{Dual viewpoint.}
Methods $W^+=W-\eta\nabla k(\nabla F(W))$ are dual-space preconditioned gradient descent \cite{maddison2021dual}, which
its authors read as a Bregman method in gradient space for smooth $k$ and which excludes normalized steepest descent;
see also \cite{kim2023mirror,chen2024lion,chen2025muon,ghane2026dual}. We use the nonsmooth case $k=\nnuc\cdot$. \emph{Exact
dynamics.} On matrix quadratics, \cite{shulgin2026quadratic} gives the one-step identities, the commuting-case map
$s_i\mapsto\abs{s_i-\bar s}$, its two-cluster invariant configurations and a complete solution in dimension two;
\cite{vasudeva2025muon,gonon2026insights,ma2026preconditioning,wang2026why,jawanpuria2026constant} contain the sign
recursions, the soft-thresholding flow and the floors of constant-length steps; finite-time convergence of
norm-constrained flows is in \cite{leplat2025norm,cortes2006finite,ma2021qualitative}, and extinction times of gradient
flows of one-homogeneous functionals in \cite{bungert2021nonlinear}. \emph{Noise.} The gradient of a randomly smoothed
support function is the expected maximizer \cite{duchi2012randomized,abernethy2014online,duboistaine2022fast}; for the
nuclear norm and Gaussian perturbations see \cite{jiang2026adaptive}. That gradient noise smooths the polar map, and
that Muon at low signal-to-noise ratio is a gradient step of size $\eta\,\E\nnuc Z/(Kd\sigma)$, are in
\cite{xie2026muon,chen2026when,jacobs2026parallel,paquette2026phases}. That sign methods drive a median, not the mean,
of the stochastic gradient to zero is known \cite{chen2020distributed}, with success-probability and symmetric-noise
conditions in \cite{bernstein2018signsgd,bernstein2019majority,safaryan2021stochastic} and counterexamples in
\cite{karimireddy2019error}; normalized steps on a location problem estimate the geometric median
\cite{cardot2013efficient} (medians in Banach spaces: \cite{minsker2015geometric}; medians of gradients:
\cite{schaipp2025tracking}); in dimension one, non-convergence of Muon for data with a nonzero third central moment is
proved in \cite{do2026muon}; variance reduction inside sign or steepest-descent steps appears in
\cite{chzhen2023signsvrg,li2026implicit}.

\section{Setting and the dual picture}\label{sec:setting}

$W\in\R^{K\times d}$, $r=\min(K,d)$, $\inner AB=\tr(A^\top B)$; $\nF\cdot$, $\nop\cdot$, $\nnuc\cdot$ are the Frobenius,
spectral and nuclear norms. For $M=U\diag(s)V^\top$ (reduced on the nonzero singular values), $\polar(M)=UV^\top$ and
$\abs M=(M^\top M)^{1/2}$; $\srank(M)=\nF M^2/\nop M^2$ and $\nr(M)=\nnuc M^2/\nF M^2$ are the stable and nuclear ranks.
$\psi(G)=\nnuc G$ is convex and $\polar(G)\in\partial\psi(G)$. A \emph{polar step of length $\eta$} is
$W^+=W-\eta\,\polar(G)$, $G=\nabla F(W)$; the \emph{nuclear-scaled step} is $W^+=W-\frac{\nnuc G}{L}\polar(G)$
\cite{carlson2015stochastic,lau2025polargrad}. $F$ is $C^2$ with minimizer $W_\ast$, and for a norm $N$ we let
$\mu_N,L_N$ be the best constants in
$\mu_NN(\Delta)^2\le D^2F(W)[\Delta,\Delta]\le L_NN(\Delta)^2$ and $\kappa_N=L_N/\mu_N$; subscripts ``op'' and $F$ refer
to the spectral and Frobenius norms. Two models are used: (LS) $F(W)=F_\ast+\frac12\tr((W-W_\ast)\Sigma(W-W_\ast)^\top)$,
$\Sigma\succ0$, i.e.\ multi-output least squares, with $E=W-W_\ast$ and $G=E\Sigma$; (GLM)
$F(W)=\frac1n\sum_i\ell_i(Wx_i)+\frac\lambda2\nF W^2$ with convex $C^2$ losses,
$c_i=\sup_z\lmax(\nabla^2\ell_i(z))<\infty$, and $f_i(W)=\ell_i(Wx_i)+\frac\lambda2\nF W^2$.

\begin{lemma}[dual picture]\label{lem:dual}
Assume $0<\muop\le\Lop<\infty$ and let $F^\ast$ be the conjugate of $F$, $D_\Phi$ the Bregman divergence of $\Phi$.
(a) A polar step is one step of mirror subgradient descent on $\psi$ in the space of gradients with mirror map
$F^\ast$: $\nabla F^\ast(G^+)=\nabla F^\ast(G)-\eta O$, $O\in\partial\psi(G)$.
(b) $D_{F^\ast}(\nabla F(W'),\nabla F(W))=D_F(W,W')$; in particular $D_{F^\ast}(0,G)=F(W)-F_\ast$.
(c) $2\muop(F(W)-F_\ast)\le\nnuc{\nabla F(W)}^2\le2\Lop(F(W)-F_\ast)$.
\end{lemma}
Item (b) is Lemma~3.2 of \cite{maddison2021dual}, item (c) the dual-norm Polyak--\L{}ojasiewicz inequality and its
converse \cite{balles2020geometry,leplat2025norm}; (a) is the nonsmooth instance of the reading of
\cite{maddison2021dual}. Two classical rules follow at once. The nuclear-scaled step with $L=\Lop$ is the mirror
Polyak step \cite{dorazio2023stochastic} for $\psi$ (whose optimal value $0$ is known), with the known rate
$1-1/\kop$. And the geometric schedule $\eta_t=\eta_0(1-1/\kop)^{t/2}$, $\eta_0=\sqrt{2\muop\bar D_0}/\Lop$ with any
$\bar D_0\ge F(W_0)-F_\ast$, gives $F(W_t)-F_\ast\le\bar D_0(1-1/\kop)^t$ when $r\ge2$; it requires $\muop$, $\Lop$ and
a bound on the initial gap. This is the geometric step-size rule of subgradient methods on sharp functions
\cite{goffin1977convergence}, transported to the dual (for spectral descent on sharp nonsmooth objectives see
\cite{yang2026convergence}).

\section{Conditioning: what no norm can buy}\label{sec:kappa}

\begin{theorem}[no free lunch, for every norm]\label{thm:nfl}
Consider (GLM) and assume that there is a sequence $W_t$ with $\max_i\nop{\nabla^2\ell_i(W_tx_i)}\to0$ (logistic and
softmax losses: $W_t=tW_1$ for generic $W_1$). Let $N$ be any norm on $\R^{K\times d}$ and $c_N=\max_{\nF\Delta=1}N(\Delta)$.
Then $\mu_N(F)=\lambda/c_N^2$ and $L_N(g)\ge L_F(g)/c_N^2$ for every convex $C^2$ function $g$. Consequently
\[
\kappa_N(F)\ \ge\ \kF(F),\qquad \frac{L_N(f_i)}{\mu_N(F)}\ \ge\ \frac{L_F(f_i)}{\mu_F(F)}=\frac{c_i\norm{x_i}^2+\lambda}{\lambda}\quad\text{for every }i .
\]
\end{theorem}
\begin{proof}
For all $W,\Delta$: $D^2F(W)[\Delta,\Delta]\ge\lambda\nF\Delta^2\ge(\lambda/c_N^2)N(\Delta)^2$. At a maximizer $\Delta_0$
of $N/\nF\cdot$, $D^2F(W_t)[\Delta_0,\Delta_0]\le\max_i\nop{\nabla^2\ell_i(W_tx_i)}\frac1n\sum_i\norm{\Delta_0x_i}^2+\lambda\nF{\Delta_0}^2
\to\lambda\nF{\Delta_0}^2$, so $\mu_N(F)=\lambda/c_N^2$. For convex $g$,
$D^2g[\Delta,\Delta]/N(\Delta)^2\ge D^2g[\Delta,\Delta]/(c_N^2\nF\Delta^2)$; take suprema.
\end{proof}

The statement covers the spectral norm ($c_N=1$: $\muop=\mu_F=\lambda$ while $\Lop(f_i)=c_i\norm{x_i}^2+\lambda r$), its
preconditioned versions $\nop{S^{1/2}\Delta P^{1/2}}$, the $\ell_\infty$ norm of sign descent, and Euclidean or
Kronecker-factored metrics; for $K=1$ and Euclidean metrics it is the no-free-lunch result of our main project. The
mechanism is that the curvature floor of a ridge-regularized model is isotropic, while a norm can only trade smoothness
against strong convexity. Any advantage of a non-Euclidean step on such problems is therefore a property of the
trajectory, not of the constants. The hypothesis on the losses cannot be dropped: for ridge least squares the norm
$\nF{\Delta\Sigma^{1/2}}$ has condition number one. The literature checks covered the spectral norm and its
preconditioned versions, not the statement for an arbitrary norm.

In the spectral geometry there is in addition a dimension barrier. For a symmetric positive definite operator $H$ on
$\R^{K\times d}$, with $q(\Delta)=\inner\Delta{H\Delta}$, one has $\Lop(H)=\max_{Q\in\cO}q(Q)$, where $\cO$ is the set of
matrices with all singular values equal to one, and $1/\muop(H)=\max_{\norm u=\norm v=1}\inner{uv^\top}{H^{-1}uv^\top}$.
The first formula is the smoothness constant for a norm \cite{balles2020geometry,islamov2026noneuclidean} reduced to the
extreme points of the spectral ball; the second is the duality between strong convexity for a norm and smoothness of
the conjugate for the dual norm \cite{kakade2012regularization}, applied to $q$. The first maximum contains the
maximization of a convex quadratic over the cube, which is NP-hard in general
\cite{islamov2026noneuclidean,balles2020geometry}.

\begin{proposition}[frame bound, rigidity, comparison]\label{prop:frame}
(a) For orthonormal families $(u_k)_{k\le r}$ in $\R^K$ and $(v_k)_{k\le r}$ in $\R^d$, with $h_k=q(u_kv_k^\top)$:
$\Lop\ge\sum_kh_k$ and $\muop\le\min_kh_k$, so $\kop\ge\sum_kh_k/\min_kh_k\ge r$.
(b) $\kop(H)=r$ if and only if $H$ is a multiple of the identity; hence a $C^2$ function has $\kop(F)=r$ iff
$F=F_\ast+\frac\mu2\nF{W-W_\ast}^2$.
(c) For every eigenpair $(\lambda,V)$ of $H$, $\Lop\ge\lambda\,\nr(V)$ and $\muop\le\lambda\,\srank(V)$; hence
$\max\{r,\ \kF/r,\ \kF\,\nr(V_{\max})/\srank(V_{\min})\}\le\kop\le r\kF$, and the lower bound $\kF/r$ is approached.
(d) For (LS) with $K\ge d$, $\Sigma=\hat\Sigma+\lambda I$, $\hat\Sigma=\frac1n\sum_ix_ix_i^\top$, and
$f_i(W)=\frac12\norm{Wx_i-y_i}^2+\frac\lambda2\nF W^2$:
$\Lop(F)=\tr\Sigma=\frac1n\sum_i\Lop(f_i)$, whereas $L_F(F)=\lmax(\Sigma)\le\frac1n\sum_iL_F(f_i)$.
\end{proposition}
\begin{proof}[Sketch]
(a) With $Q_\varepsilon=\sum_k\varepsilon_ku_kv_k^\top\in\cO$ and uniform signs, $\E q(Q_\varepsilon)=\sum_kh_k$. (b) Equality
forces $h_k=\muop$ for every frame and the vanishing of all cross terms $\inner{u_kv_k^\top}{Hu_lv_l^\top}$; polarization
gives $H=\muop\mathrm{Id}$. (c) $\lambda\nnuc V=\inner{\polar V}{HV}\le\sqrt{q(\polar V)q(V)}\le\sqrt{\Lop\lambda}\,\nF V$.
On $r\times r$ matrices, $q(\Delta)=m\nF\Delta^2+(1-m)\Delta_{12}^2-(m-b)(\tr\Delta)^2/r$ has $\kF=1/b$, $\Lop=1+m(r-1)$,
$\muop=rmb/(m+(r-1)b)$, and $\kop/\kF\to1/r$ as $b/m\to0$, $m\to0$. (d) An isometry excites all samples at once.
\end{proof}
The barrier $\kop\ge r$ is the spectral form of a folklore fact on $\ell_\infty$ \cite{sidford2018coordinate}, and
$\Lop(F)=\tr\Sigma$ is known \cite{davis2025spectral,jawanpuria2026constant,shen2025convergence}. The frame bound, the
comparison (c) and the reading of (d) below were not found; the rigidity (b) and the example showing that $\kF/r$ is
approached were added after the literature checks and were not searched for. The reading of (d) is that in the spectral
geometry least squares has \emph{no averaging gain}: the smoothness of the average is the average smoothness, so
$\kop(F)=\bar L_{\mathrm{op}}/\mu$ with $\bar L_{\mathrm{op}}=\frac1n\sum_i\Lop(f_i)=\tr\hat\Sigma+\lambda d$ and
$\mu=\lmin(\Sigma)$. This is the constant $\bar L_F/\mu$, $\bar L_F=\frac1n\sum_iL_F(f_i)=\tr\hat\Sigma+\lambda$, of
variance-reduced methods with importance sampling \cite{xiao2014proximal}, plus $(d-1)\lambda/\mu$: without ridge, the
full-batch nuclear-scaled step has the condition number of an importance-sampled variance-reduced method while paying
$n$ gradients per step. For $K<d$ the averaging gain reappears.

\section{Gradient noise: convergence to a nuclear median}\label{sec:noise}

Let $\xi$ be a random $K\times d$ matrix and
\[
\psi_\xi(G)=\E\big[\nnuc{G+\xi}-\nnuc\xi\big],\qquad \Phi_\xi(G)=\E\,\polar(G+\xi).
\]
No moment is required. $\psi_\xi$ is convex and $1$-Lipschitz for the nuclear norm, $\Phi_\xi(G)\in\partial\psi_\xi(G)$,
with $\Phi_\xi(G)=\nabla\psi_\xi(G)$ when $G+\xi$ has full rank almost surely: this is the randomized-smoothing identity
\cite{duchi2012randomized,abernethy2014online,duboistaine2022fast,jiang2026adaptive} for the spectral ball. We call
$\Med(\xi)=\arg\min_C\E[\nnuc{\xi-C}-\nnuc\xi]$ the set of \emph{nuclear medians} of $\xi$ (the median of a law in a
normed space \cite{minsker2015geometric}); $\arg\min\psi_\xi=-\Med(\xi)$ is nonempty, convex and compact, since
$\psi_\xi(G)\ge(2\Prob(\nnuc\xi\le R)-1)\nnuc G-2R$ for any $R$ with $\Prob(\nnuc\xi\le R)>\frac12$.

\begin{theorem}[stochastic polar steps are stochastic mirror descent]\label{thm:smd}
Let $F$ satisfy $0<\muop\le\Lop<\infty$, let $\xi_t$ be i.i.d.\ copies of $\xi$ independent of the past, and
$W_{t+1}=W_t-\eta_t\,\polar(\nabla F(W_t)+\xi_t)$, $G_t=\nabla F(W_t)$. For every $G^\dagger\in\arg\min\psi_\xi$, with
$W^\dagger=\nabla F^\ast(G^\dagger)$:
\begin{enumerate}[label=(\alph*),itemsep=1pt,topsep=2pt]
\item $\E[D_F(W_{t+1},W^\dagger)\mid W_t]\le D_F(W_t,W^\dagger)-\eta_t(\psi_\xi(G_t)-\min\psi_\xi)+\frac{\Lop}2\eta_t^2$;
\item $\E\,\psi_\xi(\bar G_T)-\min\psi_\xi\le\big(D_F(W_0,W^\dagger)+\frac{\Lop}2\sum_{t<T}\eta_t^2\big)/\sum_{t<T}\eta_t$ for
the weighted average $\bar G_T=\sum_{t<T}\eta_tG_t/\sum_{t<T}\eta_t$, whatever the lengths;
\item if $\sum_t\eta_t=\infty$ and $\sum_t\eta_t^2<\infty$, $G_t$ converges almost surely to a point of $-\Med(\xi)$;
\item hence, for such lengths, $W_t\to W_\ast$ if $\Med(\xi)=\{0\}$, and if $0\notin\Med(\xi)$ then almost surely
$\lim_tF(W_t)-F_\ast\ge\min_{C\in\Med(\xi)}\nnuc C^2/(2\Lop)>0$.
\end{enumerate}
\end{theorem}
\begin{proof}
With $O_t=\polar(G_t+\xi_t)$, the three-point identity of mirror descent with mirror map $F^\ast$ gives
$D_{F^\ast}(G^\dagger,G_{t+1})=D_{F^\ast}(G^\dagger,G_t)-\eta_t\inner{O_t}{G_t-G^\dagger}+D_{F^\ast}(G_t,G_{t+1})$. By
Lemma~\ref{lem:dual}(b) the first two divergences are $D_F(W_{t+1},W^\dagger)$ and $D_F(W_t,W^\dagger)$, and the last one
is $D_F(W_{t+1},W_t)\le\frac{\Lop}2\nop{W_{t+1}-W_t}^2\le\frac\Lop2\eta_t^2$. Conditionally on $W_t$,
$\E\inner{O_t}{G_t-G^\dagger}=\inner{\Phi_\xi(G_t)}{G_t-G^\dagger}\ge\psi_\xi(G_t)-\psi_\xi(G^\dagger)$. This is (a); (b)
follows by summation and convexity. (c) By the Robbins--Siegmund lemma, almost surely $D_F(W_t,W')$ converges,
simultaneously for a countable dense set of minimizers $\nabla F(W')$, and
$\sum_t\eta_t(\psi_\xi(G_t)-\min\psi_\xi)<\infty$; the iterates are bounded since
$D_F(W_t,W^\dagger)\ge\frac{\muop}2\nop{W_t-W^\dagger}^2$, a cluster point of $G_t$ is a minimizer, and the whole
sequence converges to it. (d) follows from Lemma~\ref{lem:dual}(c).
\end{proof}

The steps being bounded, no variance assumption is needed. The theorem replaces the usual condition $\E\xi=0$ by
$\Med(\xi)=\{0\}$; if $0$ is a nuclear median but not the only one, $W_\ast$ is only one of the possible limits. The
next statement makes the comparison with the mean field $W\mapsto\Phi_\xi(\nabla F(W))$ precise and bounds the offset.

\begin{proposition}[medians, zeros of the mean field, size]\label{prop:median}
(a) Every zero of $\Phi_\xi$ belongs to $-\Med(\xi)$; the converse holds at points $G$ where $G+\xi$ has full rank
almost surely, and can fail otherwise (for $\xi=\pm\frac12\diag(1,-1)$, $G=\frac12I$ minimizes $\psi_\xi$ and
$\Phi_\xi(G)=\frac12I$).
(b) If $\xi$ and $-\xi$ have the same law, $0\in\Med(\xi)$, and $\Med(\xi)=\{0\}$ if moreover $0$ is in the support of
$\xi$.
(c) If $\E\xi=0$, every $C\in\Med(\xi)$ satisfies $\nnuc C\le\E\nnuc{\xi-C}\le\E\nnuc\xi$ (the constant $1$ is sharp), so
the limit of Theorem~\ref{thm:smd} has $F(W^\dagger)-F_\ast\le(\E\nnuc\xi)^2/(2\muop)$.
(d) For a SAGA or SVRG estimator, the stochastic gradient at $(W_\ast,\text{exact table})$ is $0$ for every index: a
Dirac mass, whose median is its mean. Such estimators have a fixed point without offset, and it is the only one.
\end{proposition}
\begin{proof}
(a) $\Phi_\xi(G)\in\partial\psi_\xi(G)$, a singleton under the rank condition. (b) $\psi_\xi$ is even and convex, and
$\psi_\xi(G)-\psi_\xi(0)=\frac12\E[\nnuc{\xi+G}+\nnuc{\xi-G}-2\nnuc\xi]$ has a nonnegative integrand, positive on
$\{\nnuc\xi<\nnuc G\}$. (c) $\nnuc C=\nnuc{\E(C-\xi)}\le\E\nnuc{\xi-C}\le\E\nnuc\xi$ by Jensen's inequality and the
minimality of $C$. (d) Direct substitution.
\end{proof}
Item (b) is the matrix form of the symmetric-noise condition under which sign methods are known to converge
\cite{bernstein2019majority,jacobs2026parallel}. Item (d) identifies a fixed point and says nothing about its
attraction.

When the law of the stochastic gradient depends on $W$ (minibatches of a finite sum), Theorem~\ref{thm:smd} does not
apply as such, but the same mechanism can be followed by hand in the simplest case.

\begin{proposition}[batch one on least squares]\label{prop:lad}
Let $g_j(W)=(Wx_j-y_j)x_j^\top$ (no ridge term), $x_j\ne0$, $j$ uniform in $\{1,\dots,n\}$, and
$\Lambda(W)=\frac1n\sum_i\norm{Wx_i-y_i}/\norm{x_i}$. Then $\polar(g_j(W))\in\partial_W(\norm{Wx_j-y_j}/\norm{x_j})$, with
Frobenius norm at most one: the batch-one polar step is the stochastic subgradient method on $\Lambda$. With
$\sum\eta_t=\infty$, $\sum\eta_t^2<\infty$ it converges almost surely to a minimizer of the weighted least-distances
objective $\Lambda$ (least absolute deviations for $K=1$), which in general is not the least-squares solution.
\end{proposition}
\begin{proof}
For $r=Wx-y\ne0$, $\polar(rx^\top)=rx^\top/(\norm r\norm x)=\nabla_W\norm{Wx-y}/\norm x$
\cite[Lemma~H.1]{li2026implicit}; for $r=0$, $\polar(0)=0\in\partial_W(\norm{Wx-y}/\norm x)$. Hence
$\E_j\polar(g_j(W))\in\partial\Lambda(W)$ for every $W$, and the rest is the standard analysis of the stochastic
subgradient method in the Frobenius geometry.
\end{proof}

For $K=1$ the step is the normalized sign-error algorithm of adaptive filtering, classically associated with the
minimization of a mean absolute error. For coordinatewise sign descent with batch size one, non-convergence to the
minimizer for every step-size sequence is proved in \cite{karimireddy2019error} (Theorem~I) when all data points share
a sign pattern, with a least-squares example. For $K\ge2$ the statement was not found.

\begin{proposition}[size of the offset for minibatches; formal]\label{prop:offset}
Let $\xi_b$ be the mean of $b$ i.i.d.\ centred matrices with i.i.d.\ entries of variance $\sigma^2$ and third cumulant
$\kappa_3$, with a non-lattice law. A one-term Edgeworth expansion gives, for the minimizer $G_b^\dagger$ of $\psi_{\xi_b}$,
\begin{equation}\label{eq:offset}
G_b^\dagger=\frac{\kappa_3}{6\sigma^2b}\cdot\frac{a_{K,d}}{c_{K,d}}\;\ones\ones^\top+o\Big(\frac1b\Big),\qquad
a_{K,d}=-\E\big[\polar(Z)_{11}(Z_{11}^3-3Z_{11})\big],\quad c_{K,d}=\frac{\E\nnuc Z}{Kd},
\end{equation}
with $Z$ a standard Gaussian $K\times d$ matrix: in every entry, the scalar offset $\kappa_3/(6\sigma^2b)$ of the median
of a mean (Cornish--Fisher \cite{cornish1938moments}; $a_{1,1}=c_{1,1}=\sqrt{2/\pi}$) times $a_{K,d}/c_{K,d}$. The
objective gap at the limit is of order $1/b^2$. For a plain exponential average with parameter $\beta$ in place of the
batch mean (scalar case, frozen iterate), $1/b$ is replaced by $(1-\beta^2)/(1+\beta+\beta^2)$, which exceeds the factor
$(1-\beta)/(1+\beta)$ by which the average reduces the variance (effective batch size $b(1+\beta)/(1-\beta)$
\cite{paquette2026phases,xie2026muon}) by up to $4/3$.
\end{proposition}
\begin{proof}[Derivation]
$\polar$ is homogeneous of degree $0$, so $\Phi_{\xi_b}(0)=\E\polar(\sqrt b\,\xi_b)$; the density of $\sqrt b\,\xi_b$ is
the Gaussian one times $1+b^{-1/2}\frac{\kappa_3}{6\sigma^3}\sum_{ij}\mathrm{He}_3(z_{ij}/\sigma)+O(1/b)$ and the even terms
integrate to zero against the odd map $\polar$: $\Phi_{\xi_b}(0)=-b^{-1/2}\frac{\kappa_3}{6\sigma^3}a_{K,d}\ones\ones^\top+O(b^{-3/2})$
(by sign and permutation symmetries only the diagonal terms $\E[\polar(Z)_{ij}\mathrm{He}_3(Z_{ij})]$ survive). The
derivative of $\Phi_{\xi_b}$ at $0$ is $\sqrt b\,(c_{K,d}/\sigma)\mathrm{Id}$ to leading order
\cite{xie2026muon,chen2026when}. Solve $\Phi_{\xi_b}(G)=0$ to first order. For the average, use its cumulants
$\kappa_n(1-\beta)^n/(1-\beta^n)$ (the factors $1/(1-\beta^2)$ and $1/(1-\beta^3)$ appear in \cite{do2026muon} for
another functional of the momentum).
\end{proof}

The expansion is formal (an Edgeworth expansion against a discontinuous map) but its predictions are parameter-free and
are confirmed numerically in Section~\ref{sec:num}. Two readings. First, \emph{what matters is the skewness of the
gradient noise}: in a regression with symmetric inputs independent of the residuals, the per-sample gradient
$\varepsilon x^\top$ is symmetric in law whatever the skewness of $\varepsilon$, and the population offset is zero (on a
sample it is of the order of the sampling noise); an intercept is enough to create one. Second,
$a_{K,d}/c_{K,d}$ is of order $1/d$: in Table~\ref{tab:ac}, $d\,a_{K,d}/c_{K,d}$ lies between $1$ and $2.8$ and is about
$2$ for square matrices, and for $K=1$ one computes $a_{1,d}/c_{1,d}=3/(d+2)$ exactly. Per entry, the offset of the polar
step is thus smaller than that of sign descent, $\kappa_3/(6\sigma^2b)$ in every coordinate, by a factor of order $d$,
and its nuclear norm is at most about $2.3$ times the scalar offset for the sizes tried (smaller for wide matrices).
Orthogonalization averages the skewness. In dimension one and for the Newton--Schulz map with momentum,
non-convergence for bounded data with a nonzero third \emph{central} moment, for almost all batch sizes, is proved in
\cite{do2026muon} (Theorem~1.3), where the limits are characterized as zeros of the mean field (Proposition~6.2).

\section{Exact dynamics on least squares: complements}\label{sec:ls}

For (LS) with $K\ge d$ the polar dynamics are largely understood \cite{shulgin2026quadratic,vasudeva2025muon}: at
full-rank iterates the curvature along the polar factor is $\tr\Sigma=\Lop$ and the nuclear-scaled step is an exact
line search; if $E_0^\top E_0$ commutes with $\Sigma$ the singular values of the error follow sign recursions that do
not involve $\Sigma$, the flow is the singular-value soft-thresholding of $E_0$, and, while the amplitudes $a$ of the
error are positive, the nuclear-scaled step maps them to $\abs{a-\E_p[a]}$, the mean being taken under
$p_k=\lambda_k(\Sigma)/\tr\Sigma$, with one-step factor $\Var_p(a)/\E_p[a^2]$; the two-valued configurations
$a=\delta((1-\pi)\ones_A+\pi\ones_{A^c})$, $\pi=p(A)<\frac12$, are invariant with factor $(1-2\pi)^2$
\cite{shulgin2026quadratic,wang2026why}; the tail is slow because the step recreates tiny singular values. By the
\emph{polar flow} we mean the solution of $\dot W\in-\partial\psi(\nabla F(W))$. For (LS) it is unique: in the variable
$G=E\Sigma$ the inclusion reads $\dot G\in-\partial\psi(G)\Sigma$, the gradient flow of the convex function $\psi$ for the
inner product $\tr(A\Sigma^{-1}B^\top)$ \cite{brezis1973operateurs,bungert2021nonlinear}; since
$\polar(G)\in\partial\psi(G)$, every absolutely continuous curve with $\dot E=-\polar(E\Sigma)$ for almost every $t$ is
this solution. We add the following complements.

\begin{proposition}[rotation, conservation law, non-monotonicity]\label{prop:Q}
Let $E$ have rank $d$. (a) $\polar(E\Sigma)=\polar(E)\,Q$ with $Q=\polar(\abs E\Sigma)\in SO(d)$, and $Q=I$ iff $\abs E$
commutes with $\Sigma$. Along the flow $\dot E=-\polar(E\Sigma)$: $\frac d{dt}\frac12\nF E^2=-\tr(\Sigma^{-1}\abs G)<0$ and
each simple singular value $s_i$ of $E$, with right singular vector $w_i$, has $\dot s_i=-w_i^\top Qw_i$.
(b) If $K=d=2$, every polar step $E^+=E-\eta\polar(E\Sigma)$, of any length, conserves $s_1(E)-s_2(E)$ as long as
$\det E$ keeps its sign.
(c) If $d\ge3$, $\nop{E}$ can increase along the flow while $F$ and $\nF E$ decrease: for $E_0=\diag(300,200,1)$ and
$\Sigma=\left(\begin{smallmatrix}2&4&12\\4&9&20\\12&20&89\end{smallmatrix}\right)$, $Q_{11}=-0.0143$, and $\nop{E(t)}$
rises from $300$ to $303.27$ before decreasing. This needs an ill-conditioned error: for $d=3$,
$\frac d{dt}\nop E<0$ for every $\Sigma$ as soon as $s_1^2<(2+\sqrt3)^3s_2s_3\approx52\,s_2s_3$.
\end{proposition}
\begin{proof}
(a) $\polar(E)Q=E\abs E^{-1}\abs E\Sigma\abs G^{-1}=G\abs G^{-1}$ and $\det(\abs E\Sigma)>0$;
$\inner E{\polar(E)Q}=\tr(\abs EQ)=\tr(\Sigma^{-1}\abs G)$. (b) Split a $2\times2$ matrix into its components $M_c$, $M_a$
in the planes spanned by rotations and by reflections; its singular values are $m_c+m_a$ and $\abs{m_c-m_a}$ with
$m_\bullet=\nF{M_\bullet}/\sqrt2$, and $\det M=m_c^2-m_a^2$. If $\det E>0$ the polar factor of $E\Sigma$ is a rotation, so
the step leaves $E_a$ unchanged and $s_1-s_2=2m_a$; symmetrically if $\det E<0$. (c) Direct computation; the sign of
$Q_{11}$ was certified in exact rational arithmetic. For the last claim, $B=Q^\top\abs E=\abs G\Sigma^{-1}$ has positive
eigenvalues with product $s_1s_2s_3=:g^3$, so its trace $\sum_is_iq_i$ and the sum
$s_2s_3q_1+s_1s_3q_2+s_1s_2q_3$ of its principal $2\times2$ minors ($q_i=w_i^\top Qw_i$) are at least $3g$ and $3g^2$;
with $q_2+q_3\le1+q_1$ for a rotation this gives $q_1(s_1+s_2)(s_1+s_3)\ge3g(s_1+g)-s_1(s_2+s_3)$, which is positive
under the stated condition.
\end{proof}
The ``if'' part of $Q=I$ is in \cite{shulgin2026quadratic} (in the commuting case the curvature drops out of the
direction), and a formula of the same type for singular values along a projected polar flow is in
\cite{beneventano2026spectral}. Item (b) explains the $45^\circ$ trajectories of singular values observed for $2\times2$
problems in \cite{shulgin2026quadratic}; an $\ell_\infty$ analogue of (c) for a ball-proximal method is in
\cite{gruntkowska2025broximal}. In our searches a negative $Q_{11}$ also required a condition number of $\Sigma$ of
several hundreds.

\begin{proposition}[contraction after the first step]\label{prop:mad}
In the commuting case with positive amplitudes, the nuclear-scaled step with $L=\tr\Sigma$ satisfies
$\frac{F(W_{t+1})-F_\ast}{F(W_t)-F_\ast}=\frac{\Var_p(a(t))}{\E_p[a(t)^2]}=1-\frac{\MAD_p(a(t-1))^2}{\Var_p(a(t-1))}$ for
$t\ge1$ (the first equality is the published one-step identity), and this is at most
$1-\frac{4p_{(1)}p_{(2)}}{p_{(1)}+p_{(2)}}\le1-\frac2{\kop}$, where $p_{(1)}\le p_{(2)}$ are
the two smallest weights and $\MAD_p$ is the mean absolute deviation. The constant is sharp: it is a supremum, attained
iff $d=2$, where the factor is the exact $((\kappa-1)/(\kappa+1))^2$, $\kappa=\lambda_1/\lambda_2$, proved in
\cite{shulgin2026quadratic} for every $t\ge1$ and every full-rank start. The same bound holds for the second step from
every full-rank start with $W_1\ne W_\ast$, commuting or not (new for $d\ge3$). At $t=0$ the factor can be arbitrarily
close to the certified $1-1/\kop$ \cite{shulgin2026quadratic}. Conversely every invariant ray of
$a\mapsto\abs{a-\E_pa}$ is two-valued, and for $d\ge3$ these rays are linearly unstable.
\end{proposition}
\begin{proof}[Sketch]
After one step $a(t)=\abs X$ with $X$ centred under $p$. With $S=\E_p[X_+]=\E_p[X_-]$, convexity gives
$\E_p[X_\pm^2]\le S^2/\min p_k$ on each side, the two minima being attained at distinct indices, so
$(\E_p\abs X)^2/\E_pX^2\ge4p_{(1)}p_{(2)}/(p_{(1)}+p_{(2)})$. Without commutation, after one line-search step
$G_1=\polar(G_0)M_1$ with $M_1$ symmetric of trace zero \cite{shulgin2026quadratic}; the same argument applied to the
positive and negative parts of $M_1$, with Ky Fan's inequality, gives the bound for the second step. The derivative of
the map on a ray is $\diag(\pm1)(I-\ones p^\top)$, with eigenvalues $\pm1$ besides $1-2\pi$ and $0$.
\end{proof}

\begin{proposition}[extinction time of the polar flow]\label{prop:ext}
For a $C^2$ function with $0<\muop\le\Lop<\infty$, a solution of $\dot W\in-\partial\psi(\nabla F(W))$ reaches $W_\ast$ at a finite time
$T$ with $\nop{W_0-W_\ast}\le T\le\sqrt{2(F(W_0)-F_\ast)/\muop}\le\sqrt{\kop}\,\nop{W_0-W_\ast}$. For (LS) the lower bound
is attained in the commuting case, whatever $\Sigma$; from a full-rank start the flow is the segment
$E_0-t\,\polar(G_0)$ until the first rank drop, at a time
$\tau_\ast\in\big[s_d(E_0),(\prod_is_i(E_0))^{1/d}\big]$; for $d=2$, $T\le\frac{9\sqrt2}8\nop{E_0}$ whatever $\Sigma$; and
no bound $T\le C\nop{E_0}$ with $C$ independent of both $d$ and $\Sigma$ holds.
\end{proposition}
\begin{proof}
$\frac d{dt}(F-F_\ast)=-\nnuc G\le-\sqrt{2\muop(F-F_\ast)}$ \cite{leplat2025norm} and $\nop{\dot W}\le1$. Along
$E_0-t\,\polar(G_0)$ the gradient is $\polar(G_0)(\abs{G_0}-t\Sigma)$, with the same polar factor while
$\abs{G_0}-t\Sigma\succ0$ (the rank-drop times are those of \cite{shulgin2026quadratic}); $\tau_\ast$ is the smallest
eigenvalue of $\Sigma^{-1/2}\abs{G_0}\Sigma^{-1/2}$, whose eigenvalues have product $\det\abs{E_0}$, and
$\abs{G_0}\succeq s_d(E_0)\Sigma$. For a rank-one error the flow is the gradient flow of $\frac12e^\top\Sigma e$
parametrized by arc length, of length at most $\sqrt d\,\norm{e_0}$ and possibly of order
$\min\{\sqrt d,\sqrt{\log\kappa(\Sigma)}\}\norm{e_0}$ \cite{gupta2021path}. For $d=2$ both singular values decrease at the
same speed (Proposition~\ref{prop:Q}), so the error has rank one and norm $s_1-s_2$ at
$\tau_\ast\le\sqrt{s_1s_2}$, and $T\le\sqrt{s_1s_2}+\sqrt2(s_1-s_2)$.
\end{proof}
The first sentence is known: for the flow of any norm, finite extinction and the bound
$\sqrt{2(F(W_0)-F_\ast)/\mu}$ under the dual-norm Polyak--\L{}ojasiewicz inequality are in \cite{leplat2025norm}, and
for quadratics the two bounds are instances of the extinction-time bounds of \cite{bungert2021nonlinear}; the equality
in the commuting case is read on the soft-thresholding flow of \cite{vasudeva2025muon}. Added here: the straight
segment of the flow until the first rank drop, the bounds on $\tau_\ast$ and the bound for $d=2$. From generic
(Gaussian, full-rank) starts, $T/\nop{E_0}\in[1,1.4]$ in all instances tried, far below $\sqrt{\kop}$ (up to $24$);
larger ratios need a nearly rank-one error (up to $\sqrt2$ for $d=2$, and close to $\sqrt d$ in the rank-one limit
with a widely spread spectrum). The head of the dynamics is thus nearly curvature-free beyond the commuting case, which
remains to be explained.

\section{Numerical illustrations}\label{sec:num}

\begin{table}[t]
\centering\small
\begin{tabular}{@{}lccccccccc@{}}
\toprule
$(K,d)$ & $(1,1)$ & $(2,2)$ & $(2,3)$ & $(4,4)$ & $(8,8)$ & $(16,16)$ & $(32,32)$ & $(4,16)$ & $(16,64)$\\
\midrule
$a_{K,d}$ & 0.798 & 0.391 & 0.286 & 0.176 & 0.0737 & 0.0285 & 0.0107 & 0.0389 & 0.0053\\
$c_{K,d}$ & 0.798 & 0.573 & 0.504 & 0.412 & 0.295 & 0.210 & 0.149 & 0.240 & 0.121\\
$a_{K,d}/c_{K,d}$ & 1 & 0.683 & 0.568 & 0.428 & 0.250 & 0.136 & 0.072 & 0.162 & 0.044\\
$d\cdot a_{K,d}/c_{K,d}$ & 1 & 1.37 & 1.70 & 1.71 & 2.00 & 2.17 & 2.29 & 2.59 & 2.80\\
\bottomrule
\end{tabular}
\caption{Constants of the first-order offset \eqref{eq:offset}. Monte Carlo, with standard errors of $a_{K,d}$ between
$1.4\cdot10^{-4}$ and $7.5\cdot10^{-4}$ (up to $3\%$ of $a_{K,d}$ for the largest sizes), except the column $(1,1)$,
which is exact, and the column $(2,3)$, which gives the exact values obtained by quadrature (Monte Carlo:
$a_{2,3}=0.2869\pm0.0006$); the column $(2,2)$ agrees with its exact values. The offset per entry, in units of the
scalar offset $\kappa_3/(6\sigma^2b)$, is $a_{K,d}/c_{K,d}$, of order $1/d$; for sign descent it is $1$.}\label{tab:ac}
\end{table}

\begin{figure}[t]
\centering
\includegraphics[width=\textwidth]{figT2_noise}
\caption{Stochastic polar steps on a finite-sum least-squares problem ($n=2000$, $K=2$, $d=3$, intercept, skewed
residuals). (a) Batch one, same decreasing lengths: with the minibatch gradient the averaged iterate converges to the
least-distances solution $W_{\mathrm{LAD}}$ of Proposition~\ref{prop:lad}; with the SAGA estimator it converges to the
least-squares solution $W_{\mathrm{LS}}$. (b) Batch size times the nuclear norm of the gradient at the limit point:
zeros of the mean field (line), stochastic runs with decreasing lengths (dots), SAGA estimator with the nuclear-scaled
length (squares; the gradient norms are at rounding level, below $3\cdot10^{-15}$). (c) Scalar case with momentum:
median of an exponential average of skewed noise, with the two factors of Proposition~\ref{prop:offset}.}\label{fig:noise}
\end{figure}

\paragraph{Offset.} For $2\times3$ matrices with entries $\mathrm{Exp}(1)-1$ ($\sigma=1$, $\kappa_3=2$),
\eqref{eq:offset} predicts $\sqrt b\,\nF{\Phi_{\xi_b}(0)}\to0.234$ and $b\,\nnuc{G_b^\dagger}\to0.464$. Computing the zero of
the mean field directly, with exactly sampled batch means and $1.5\cdot10^6$ samples, gives $0.230$--$0.237$ and
$0.450$--$0.498$ for $b=2,\dots,256$; an independent implementation with $3\cdot10^6$ samples gives $0.230$--$0.241$ and
$0.449$--$0.483$ for $b$ up to $1024$. The fixed point is proportional to the all-ones matrix (cosine above $0.999$ for
$b\le64$; $0.998$ and $0.996$ at $b=128$ and $256$, where the Monte-Carlo error of the root is no longer negligible),
and with symmetrized noise the mean field at $0$ vanishes up to the Monte-Carlo error ($10^{-3}$).

The dichotomy of Theorem~\ref{thm:smd} is visible on a finite sum (Figure~\ref{fig:noise}). For a least-squares
problem with an intercept and skewed residuals, $\nF{W_{\mathrm{LAD}}-W_{\mathrm{LS}}}=0.170$; batch-one polar steps
with decreasing lengths end (iterates averaged over the last half of the run and over $32$ runs) at distance
$2.6\cdot10^{-3}$ from the least-distances solution and $0.169$ from the least-squares solution, while the same lengths
with the SAGA estimator end at $2.1\cdot10^{-4}$ from the least-squares solution. For $b=1,2,4,8,16,32$ the zero of the
mean field (estimated from $2\times2\cdot10^5$ sampled batches for $b\ge2$, with a spread up to $8\cdot10^{-4}$) has
$\nnuc{\nabla F}$ equal to $0.168$, $0.120$, $0.069$, $0.036$, $0.019$, $0.0096$ ($b$ times this: $0.168$, $0.240$,
$0.276$, $0.285$, $0.303$, $0.308$), and the stochastic runs (same averaging, $64$ runs) end at $0.166$, $0.118$,
$0.068$, $0.035$, $0.018$, $0.0092$; with the SAGA estimator and the nuclear-scaled length the gradient norm ends below
$2.4\cdot10^{-15}$ for every $b$. In this design the per-sample gradients have rank one and the entries of the noise
are not independent, so \eqref{eq:offset} does not apply as such: what is observed is the $1/b$ regime, a plateau of
$b$ times the offset, approached for $b$ large compared with the size of the matrix.

\paragraph{Momentum.} Exact medians of the stationary exponential average of $\mathrm{Exp}(1)-1$ noise are
$-0.0466$, $-0.0227$, $-0.0112$ for $\beta=0.8,0.9,0.95$; the third-cumulant factor gives $-0.0492$, $-0.0234$, $-0.0114$
and the effective-batch heuristic $-0.0370$, $-0.0175$, $-0.0085$ (Figure~\ref{fig:noise}(c)).

\paragraph{Dynamics.} For $K=d=100$ and a spectrum $1/k$ ($\kF=100$, $\kop=519$), from a generic start, the relative
objective gap after $10$ and $100$ iterations is $4.0\cdot10^{-3}$ and $1.3\cdot10^{-4}$ for the nuclear-scaled polar step
against $0.23$ and $9.0\cdot10^{-3}$ for gradient descent with step $1/\lambda_1$; after $1000$ iterations the order is
reversed ($4.5\cdot10^{-8}$ against $1.8\cdot10^{-11}$). Over iterations $2500$ to $3000$ the exponent per step of the
polar run, $0.0081$, is within $5\%$ of that of the flattest two-valued invariant configuration,
$-2\log(1-2\lambda_d/\tr\Sigma)$, and $2.5$ times smaller than that of gradient descent, the ratio
$\tr\Sigma/(2\lambda_1)=2.6$ predicted on that configuration (the slow regime of \cite{shulgin2026quadratic}, written
with the stable rank of the data). This agreement is an observation, and only of the exponent: at the end of the run
the error is nearly of rank one ($99.9\%$ of $\nF E^2$ on its top singular value) but not aligned with the flattest
eigenvector (cosine $0.54$), and the configurations of the ten flattest directions have exponents within $10\%$ of
each other. Moreover the invariant configurations are unstable for the commuting map (Proposition~\ref{prop:mad}), and
in commuting runs the tail is governed by the direction that carries the largest amplitude after the first few
iterations (exponent close to $-2\log(1-2p_k)$), not by the flattest one.

\section{Discussion}\label{sec:discussion}

\paragraph{What this says about Muon.} On the side of \emph{speed}, the spectral geometry cannot improve a worst-case
constant on ridge-regularized classification models (Theorem~\ref{thm:nfl}), and its condition number is at least
$\min(K,d)$; what polar steps buy is a property of trajectories --- a nearly curvature-free reduction of the error in
the head of the run, paid by a tail that, on the flattest invariant configuration and in our runs, is slower than
gradient descent with step $1/\lambda_1$ when $\tr\Sigma/\lambda_1$, the stable rank of the data, exceeds two. On the
side of \emph{precision}, a polar step with stochastic gradients faces three distinct floors: a length floor for
constant lengths (quadratic in the length, removed by a nuclear-scaled or decreasing length), the fluctuations of any
constant-step stochastic method, and the median offset of Theorem~\ref{thm:smd}, which is the same for all schedules
with $\sum_t\eta_t=\infty$, $\sum_t\eta_t^2<\infty$. When the gradient noise is symmetric, $W_\ast$ is a possible limit,
and the only one if $0$ is in the support of the noise (Proposition~\ref{prop:median}(b)). Otherwise the offset is
formally of order $1/b$ when the noise has a nonzero third cumulant (Proposition~\ref{prop:offset}), and smaller per
entry than for sign descent by a factor of order $d$. In the scalar quasi-static computation of
Proposition~\ref{prop:offset}, a plain exponential average multiplies it by
$(1-\beta^2)/(1+\beta+\beta^2)\approx\frac23(1-\beta)$; with the Nesterov look-ahead of Muon implementations the
first-order factor is larger ($0.125$ and $0.052$ at $\beta=0.9$ and $0.95$, against $0.070$ and $0.034$). A finite-sum
variance-reduced estimator has a fixed point without offset. In experiments of the companion study with constant
parameters, only finite-sum estimators combined with a nuclear-scaled or clipped length reach rounding precision;
Figure~\ref{fig:noise}(b) shows the same for the limit point. This is the precise sense in which variance reduction
can improve the precision of Muon-type steps: not by reducing a variance term in a rate, but by having the minimizer
as its fixed point (the convergence to it is observed in Figure~\ref{fig:noise}, not proved).

\paragraph{Limitations.} All statements are for strongly convex objectives. Theorem~\ref{thm:smd} assumes additive
noise and no momentum; for minibatches of a finite sum the law of the noise depends on the iterate, and we only treat
batch one on least squares. Proposition~\ref{prop:offset} is a formal expansion, supported by simulations; its momentum
part is scalar, quasi-static and for a plain exponential average. No convergence of a variance-reduced polar method is
proved: Proposition~\ref{prop:median}(d) only identifies its fixed point. The scaling of $a_{K,d}/c_{K,d}$ in $1/d$ is
numerical (exact for $K=1$), and it concerns noise with independent entries in the regime $b\to\infty$ at fixed size;
minibatch gradients of rank smaller than $\min(K,d)$ are outside it. ``Not found'' refers to automated literature
checks of 12 and 16 November 2026, which opened about sixty papers; several of the closest ones are a few weeks old.
Statements added during the revision (Theorem~\ref{thm:nfl} for an arbitrary norm, the rigidity and the sharpness
example of Proposition~\ref{prop:frame}, the second-step bound and the instability of the invariant rays in
Proposition~\ref{prop:mad}, the condition for $d=3$ in Proposition~\ref{prop:Q}, the bounds on $\tau_\ast$ and the bound
for $d=2$ in Proposition~\ref{prop:ext}) were not the object of a separate search.

\paragraph{Open questions.} Whether generic trajectories of the line-search polar iteration converge to the flattest
two-valued configuration; a bound on the extinction time of the flow of order $\sqrt{\log\kappa}$ or $\sqrt d$; a proof
of \eqref{eq:offset} and of its dimension dependence; a Lyapunov function for a SAGA table driven by nuclear-scaled
polar steps; and whether the dependence of the final loss on the skewness of the minibatch gradient, on $b$ and on
$\beta$ is visible when training a network.

\paragraph{Reproducibility and use of AI tools.} Proofs of all statements, the status of each auxiliary result and the
record of the review are in the companion notes (\texttt{muon/notes\_spectral\_theory.pdf}); every number comes from the
scripts of \texttt{muon/theory\_checks/}. A large language model (Claude, Anthropic) was used, under the supervision of
the authors, to derive and check the statements, to write the scripts, and to run the adversarial review (six
independent passes re-deriving each proof and opening the cited sources, then three passes on this paper, whose
corrections are included); statement numbers quoted from recent preprints are those read during that review, and the
exact constants quoted for $K\le2$ in Table~\ref{tab:ac} and the identity $a_{1,d}/c_{1,d}=3/(d+2)$ come from it.

\bibliographystyle{plain}
\bibliography{paper_spectral_muon}

\end{document}
